\chapter{Orientation}

\section{Purpose of the Framework}

The Unified Rigidity Framework is an exposition layer for a single structural
claim pattern:

\[
\text{entropy growth}
+
\text{bounded information capacity}
+
\text{coercive structural control}
\Longrightarrow
\text{nontrivial refinement depth}.
\]

The purpose of the framework is not to declare that every problem in every
domain has been solved. Its purpose is to give a common language for comparing
systems in which a large space of possible configurations cannot be collapsed
instantaneously by local or capacity-bounded operations.

The central idea is that many hard problems have the same visible shape.
There is a large configuration space. There is a process that attempts to
simplify, refine, classify, or eliminate possibilities. There is a limit on
how much information the process can extract at one step. There is also a
rigidity mechanism preventing defects from disappearing without cost. When
these ingredients are present at the same time, the process must take depth.

This book records that shape in a form that can be audited. Each theorem or
principle is separated into hypotheses, conclusion, example, and boundary
statement. The boundary statement is part of the mathematical discipline: it
says what the result does not prove.

\section{Core Structural Objects}

The framework uses a small set of recurring objects.

\begin{definition}[State space]
A state space \(X\) is the collection of admissible configurations for the
system under study. Depending on the domain, \(X\) may be finite, measurable,
topological, algebraic, spectral, or syntactic.
\end{definition}

\begin{definition}[Refinement system]
A refinement system is a pair \((X,\mathcal R)\), where \(X\) is a state space
and \(\mathcal R\) is a class of admissible refinement maps
\[
R:X\to X.
\]
A refinement map represents one allowed step of structural update.
\end{definition}

A refinement step may be an algorithmic update, a partition refinement, a
local graph-coloring update, a normalization step, a spectral projection, or a
domain-specific simplification. URF does not identify all of these processes.
It records the common question: how much uncertainty can one admissible step
remove?

\begin{definition}[Entropy]
An entropy functional \(H:X\to \mathbb R_{\ge 0}\) measures residual
uncertainty, unresolved multiplicity, or remaining defect size.
\end{definition}

For a finite configuration family \(\Omega\), the simplest model is
\[
H(\Omega)=\log |\Omega|.
\]
For a probability distribution \(p\), the standard discrete entropy is
\[
H(p)=-\sum_x p(x)\log p(x).
\]
The framework does not require these to be the only entropy notions. It
requires that the chosen entropy be explicit and compatible with the claimed
refinement bound.

\begin{definition}[Capacity]
A capacity bound \(C\) is an upper bound on the amount of relevant information
that an admissible refinement step may extract or eliminate.
\end{definition}

Capacity is not the same as total computational power. It is a local or
operational throughput constraint. A global algorithm may be powerful, but if
a theorem assumes that each permitted update sees only bounded local data,
then the relevant per-step capacity is bounded by that admissible view.

\begin{definition}[Rigidity functional]
A rigidity functional is a map
\[
F:X\to \mathbb R_{\ge 0}
\]
whose vanishing identifies a defect-free or terminal state, and whose positive
value measures the cost of unresolved structure.
\end{definition}

A rigidity functional is useful only when it is connected to entropy,
capacity, or coercivity. A named functional without such a connection is only
notation.

\section{Capacity--Entropy Inequality}

The basic bookkeeping principle is that unresolved information cannot be
removed faster than the allowed process can access it.

\begin{definition}[Capacity--entropy gap]
Given entropy \(H\) and per-step capacity \(C\), define the gap
\[
\Phi(x)=H(x)-C.
\]
The sign and interpretation of \(\Phi\) depend on the level being measured:
single-step feasibility, cumulative feasibility, or terminal feasibility.
\end{definition}

At a single step, the relevant statement is not that all configurations must
satisfy \(H(x)\le C\). Rather, the operational statement is:

\begin{quote}
If a refinement step can remove at most \(C\) units of relevant information,
then a state with \(H(x)\) units of unresolved information cannot be fully
collapsed in fewer than \(H(x)/C\) such steps.
\end{quote}

This is the source of the depth lower bound. It is deliberately modest. It
does not rule out nonlocal algorithms, global certificates, quantum procedures,
or any process outside the stated admissible class.

\section{Locality Constraint}

Locality is the condition that an update cannot inspect the entire structure
at once.

\begin{definition}[\(R\)-local refinement]
A refinement rule is \(R\)-local if the update assigned to a location \(v\)
depends only on data contained in the radius-\(R\) neighborhood of \(v\).
\end{definition}

In graph settings, this means that a vertex update depends on a bounded
neighborhood. In logical settings, this is reflected by bounded-variable or
bounded-radius formulas. In physical settings, locality may mean finite
propagation, finite detector access, or bounded operational coupling.

Locality is not a universal law of all mathematics. It is an assumption. The
framework becomes meaningful only after the admissible locality class is named.

\section{Spectral Coercivity}

A second recurring mechanism is coercivity. In analytic or spectral settings,
coercivity states that defect has positive cost away from the kernel.

\begin{definition}[Spectral coercivity]
Let \(L\) be a self-adjoint operator on a Hilbert space, and let
\(f\perp \ker L\). A spectral coercivity estimate has the form
\[
\langle f,Lf\rangle \ge \lambda \|f\|^2
\]
for some \(\lambda>0\).
\end{definition}

The role of coercivity is to prevent defect from vanishing freely. If a
system has a positive gap, then movement away from the rigid or terminal
subspace costs energy, rank, mass, entropy, or another controlled quantity.
The name of the quantity changes by domain. The structural role is the same:
coercivity turns qualitative obstruction into a quantitative lower bound.

\section{EntropyDepth}

\begin{definition}[EntropyDepth]
Let \(x_0,x_1,\ldots\) be a refinement trajectory and let \(T\) be the target
set of terminal states. The EntropyDepth of \(x_0\) is
\[
\ED(x_0)=\inf\{t\ge 0:x_t\in T\}.
\]
\end{definition}

If each step removes at most \(C\) units of entropy and the initial unresolved
entropy is \(H(x_0)\), then the basic lower bound is
\[
\ED(x_0)\ge \frac{H(x_0)}{C}.
\]

This estimate is intentionally simple. Its value is not that it is surprising
by itself. Its value is that many different problems can be reduced to the
same checklist:

\begin{enumerate}
\item identify the configuration space;
\item define the entropy or unresolved defect;
\item prove a per-step capacity bound;
\item prove locality or admissibility of the refinement class;
\item prove coercive control of the defect;
\item conclude a depth lower bound only inside that admissible class.
\end{enumerate}

\section{The Refinement Wall}

The Refinement Wall is the name for the obstruction that appears when all
three constraints are simultaneously present.

\begin{theorem}[Refinement Wall, schematic form]
Suppose a family of systems has entropy growing at least linearly in a size
parameter \(n\). Suppose further that admissible refinement steps are local and
remove at most \(C\) units of relevant entropy per step, with \(C\) independent
of \(n\). Then any admissible refinement trajectory reaching a terminal state
has depth
\[
\ED(x_n)=\Omega(n).
\]
\end{theorem}

\begin{proof}
Let \(H(x_n)\ge an\) for some \(a>0\). If each admissible step removes at most
\(C\) units of entropy, then after \(t\) steps the total removable entropy is
at most \(tC\). Reaching a terminal state requires \(tC\ge an\). Hence
\(t\ge (a/C)n\), so \(\ED(x_n)=\Omega(n)\).
\end{proof}

The theorem is schematic because each concrete domain must supply its own
notion of entropy, admissible refinement, terminal state, and capacity bound.

\section{Reduction Principle}

The reduction principle says that a domain-specific problem may be studied by
translating it into the common URF tuple
\[
(\text{state space},\text{refinement rule},\text{entropy},
\text{capacity},\text{rigidity functional},\text{frontier status}).
\]

A successful reduction does not automatically solve the original problem.
It records exactly what has been transferred. A complete theorem requires
all hypotheses in the target domain, not merely the existence of similar
language.

\section{Scope and Limits}

The framework applies only under explicit assumptions:

\begin{itemize}
\item bounded information extraction;
\item fixed or controlled locality;
\item a valid entropy or defect functional;
\item coercive control of the relevant defect;
\item a specified terminal condition.
\end{itemize}

The framework does not claim:

\begin{itemize}
\item that all algorithms are local;
\item that all systems obey bounded capacity;
\item that all defects admit coercive control;
\item that all domain translations are complete;
\item that a frontier is solved merely because it has been named.
\end{itemize}

Escaping a URF lower bound means violating or bypassing at least one listed
hypothesis. This is not a weakness of the framework. It is the audit boundary.

\section{Organization of the Book}

The book proceeds from primitives to mechanisms.

\begin{enumerate}
\item Core invariants: entropy, capacity, depth, and spectral gap.
\item Operational rigidity: how local indistinguishability becomes a global
constraint.
\item Non-amplification: why bounded updates cannot create unlimited
information.
\item EntropyDepth: how capacity bounds imply refinement-depth lower bounds.
\item Terminal obstructions: named configurations where refinement cannot
short-circuit the wall.
\item Applications: computational, spectral, physical, and structural examples.
\item Appendices: definitions, notation, dependency graph, theorem index, and
framework summary.
\end{enumerate}

\section{Structural Thesis}

The structural thesis of URF is:

\[
\text{Entropy Growth}
+
\text{Bounded Capacity}
+
\text{Coercivity}
\Longrightarrow
\text{Depth}.
\]

Every later chapter is an instance, refinement, limitation, or boundary case
of this thesis.

\begin{remark}[Boundary]
This chapter is an orientation chapter. It does not prove unrestricted
Chronos-RR, unrestricted H4.1/FGL, P vs NP, or any Clay problem. It only fixes
the vocabulary used by later chapters.
\end{remark}
